\begin{abstract}
Multimodal language models are strong at tasks that ask what is in an image and weaker, in ways
their image-task scores do not predict, at tasks that ask about geometry. We use origami to measure
the difference. We release \textbf{CP2Seq}, a benchmark and dataset of 600 origami samples in which
a model is given a crease pattern and the final folded state and must recover the sequence of folds
between them. The task is trivial to generate and hard to solve: folding forward is one pass of an
engine that records what it did, while deciding whether a crease pattern is reachable by simple
folds is NP hard, and determining the layer ordering of a flat folding is NP hard even given a valid
mountain-valley assignment. Ground truth is constructed rather than annotated, and every sample
rebuilds byte-identically from its seed, so the corpus ships as a generator rather than as data. The
environment the model works against is also the verifier: it executes one candidate fold and either
returns the resulting state or refuses with a named reason, such as the sheet would tear. It
performs no search, so proposing, pruning and backtracking stay with the model. We provide
progressive tiers of tool assistance, and because one crease pattern admits many valid folded states
we score answers by replaying them, up to the plane isometry group. At low reasoning effort no model
we evaluated solved any sample beyond the easy stratum, and on that stratum a deterministic
breadth-first search over the same action set solves 54.7\% against the best model's 20.4\%. The
dominant failure is not an exhausted budget but revisiting states already seen, which ends 48\% of
that model's attempts.
\end{abstract}
